\thispagestyle{empty}

\begin{center}
{\small\instituicao\par}
{\small\unidadeacademica\par}
{\small\disciplina\ -- \semestre\par}

\vfill

{\bfseries\MakeUppercase{\titulorelatorio}\par}
\vspace{0.5cm}
{\subtitulorelatorio\par}

\vspace{1.2cm}
{\equipe\par}
{\autores\par}

\vspace{0.65cm}
{\small Docente: \docente\par}
\end{center}

\vfill
\begin{center}
{\small Brasília\par}
{\small 2026\par}
\end{center}

\clearpage

\thispagestyle{empty}
\tableofcontents

\clearpage
\setcounter{page}{1}
\section*{RESUMO}
\addcontentsline{toc}{section}{RESUMO}

\begin{resumorelatorio}
Este relatório registra o desenvolvimento de um produto orientado à redução de avarias em louças sanitárias no transporte entre o distribuidor e a obra. O pré-desenvolvimento delimitou a oportunidade sem fixar solução: reduzir a quebra e seus efeitos sobre reposição, prazo, desperdício e segurança, preservando a viabilidade comercial. Foram identificados os instrumentos legais e normativos que restringem o manuseio, a amarração da carga e a destinação do resíduo, e foi adotado um custo-meta provisório por ciclo de uso. No projeto informacional, a equipe adotou produção própria sob a estratégia \textit{Assemble to Order} (ATO), descreveu o ciclo de vida operacional em treze fases e converteu 44 necessidades em 23 requisitos dos clientes. Esses requisitos foram classificados pelo diagrama de Kano, ponderados pelo diagrama de Mudge e ordenados pelo gráfico de Pareto. O benchmarking comparou quatro referências já em uso ou documentadas. O desdobramento da função qualidade produziu 23 requisitos de projeto, com matrizes de relacionamento e de correlação, e encerrou a fase em especificações-meta. A louça sanitária permanece como objeto protegido, e não como o produto em desenvolvimento.

\vspace{0.5\baselineskip}
\noindent\textbf{Palavras-chave:} desenvolvimento de produtos. Louças sanitárias. Avarias. Transporte. QFD. Projeto informacional.
\end{resumorelatorio}

\vfill
\begin{center}
{\footnotesize Brasília, 2026. Versão \versaodocumento. Situação: \statusprojeto.}
\end{center}

\clearpage
